% !TEX root = ../main.tex
\chapter{Conclusion}
\label{ch:conclusion}

This thesis has studied inference for extreme expectile-based risk measures in
a distributed heavy-tailed setting. The construction deliberately stays close
to the optimal-pooling framework of \textcite{DaouiaPadoanStupfler2021}: it
first forms a geometrically pooled Weissman estimator of an extreme quantile
and only then applies the high-expectile/high-quantile bridge. Looking back,
the main value of the two-weight formulation is that it separates two roles
that are easy to conflate: \(\bm\omega\) controls the pooled-Weissman
component, while \(\bm\nu\) enters only through the Hill estimator used in the
outer bridge. This separation prevents the diagonal convention
\(\bm\nu=\bm\omega\) from being mistaken for an optimality result and makes
clear which part of the estimator is governed by the source pooling theory.

\section{Main findings}
\label{sec:conclusion:findings}

The theoretical argument is useful precisely because it is explicit about what
has to be controlled. The log error splits into the pooled-Weissman error, the
plug-in error from evaluating \(\psifn\) at a pooled Hill estimator, and the
population error in replacing the expectile by the asymptotic quantile bridge.
In the iid common-marginal distributed regime, under the extreme-Weissman
conditions and the bridge-rate condition, the two bridge terms are lower order
on the pooled-Weissman scale. The limiting law therefore depends on
\(\bm\omega\), but not on \(\bm\nu\).

This is the main message of \Cref{ch:pooled-extreme}: first-order variance and
AMSE optimisation transfer to \(\bm\omega\) from the
Daouia--Padoan--Stupfler theory, while \(\bm\nu\) is first-order unidentified
in this regime. The practical convention
\(\bm\nu_n^k=(k_1/k,\ldots,k_m/k)^\top\) should therefore be read as a
transparent implementation choice, not as an optimal bridge-Hill rule.

The thesis then extends the deterministic result to source-admissible
estimated \(\bm\omega\) weights. This transfer is deliberately conditional:
when the Daouia--Padoan--Stupfler construction supplies consistent estimated
weights and consistent plug-in centring and variance objects, the expectile
estimator inherits the same first-order studentised limit. The resulting
confidence interval is obtained by inverting the transferred log-scale
expectile statistic. It is not a direct reuse of a quantile interval formula,
and it does not introduce an expectile-specific plug-in bias or variance
estimator.

\section{Finite-sample lessons}
\label{sec:conclusion:finite-sample}

The simulation study in \Cref{ch:finite-sample} is intentionally narrow. It
uses the practical bridge-Hill convention and varies only the pooled-Weissman
weights. Within this design, Daouia--Padoan--Stupfler variance weights provide
a stable practical baseline. In the representative imbalanced grid, they make
the distributed estimator track the centralised benchmark closely for point
estimation, whereas equal pooled-Weissman weights are less efficient under
strong imbalance.

The finite-sample role of AMSE weights is more conditional. In the main
equal-threshold grid, oracle and plug-in AMSE weights are nearly
indistinguishable from variance weights. Under heterogeneous local threshold
fractions, AMSE weights can matter more, and in some rows they improve either
coverage or log RMSE. They are not uniformly superior, however, and the
diagnostics show that signed or unstable weights can appear. Thus the
simulation evidence supports reporting AMSE weights as a meaningful
source-based alternative, not as a uniformly preferable default.

The interval results are more cautionary. The first-order log-scale intervals
under-cover in several finite-sample designs, and the centralised oracle
benchmark also under-covers. The decomposition diagnostics explain why this is
not only a distributed-splitting issue. At the simulated sample sizes, the
terms \(B_n(\bm\nu)\) and \(C_n\), although asymptotically negligible under
the theorem's scaling, can still be visible. In particular, finite-sample
bridge error can affect calibration even when point estimation is competitive.
The interval study should therefore be read as a diagnostic of the plug-in
route rather than as evidence of strong finite-sample nominal coverage.

\section{Limitations and future work}
\label{sec:conclusion:future}

